% ---------------------------------------------------------------------------
% Section 3: Model. Every symbol, value and rule here is the shipped code's:
% include/agent.h (utility), src/cpp/population.cpp (sampling, calibration),
% src/cpp/simulation.cpp (swipe log), src/cpp/matching.cpp (deferred
% acceptance, theta update), src/python/ncf.py (ranker), src/python/metrics.py.
% ---------------------------------------------------------------------------

\section{Model}
\label{sec:model}

The market has $N$ agents, $N/2$ on each side: a proposing side $M$ and an
accepting side $W$. It runs in discrete \emph{rounds}. In each round a learned
ranker gives every agent a shortlist of $k$ agents from the other side, and a
truncated deferred-acceptance pass clears the market over those shortlists and
nothing else. Throughout, a \emph{round} is one pass of the market loop of
Algorithm~\ref{alg:round}, and a \emph{training pass} is one sweep of the
ranker's optimizer over one round's training log. The shipped configuration has
$N = 1000$, $k = 20$ and one training pass per round. Table~\ref{tab:params}
lists every parameter with its shipped and swept values.

\subsection{Agents and utility}
\label{sec:utility}

Agent $i$ carries four traits $x_i = (s_i, a_i, c_i, t_i)$: status,
attractiveness, sociability and age, drawn independently as
\begin{equation}
s_i = e^{0.8 Z}, \quad a_i = [0.5 + 0.15 Z]_0^1, \quad
c_i = [0.5 + 0.15 Z]_0^1, \quad t_i = 30 + 6 Z,
\label{eq:traits}
\end{equation}
with a fresh standard normal $Z$ for each draw and $[\cdot]_0^1$ clipping to the
unit interval. Each trait is standardized against the whole population, so
$z(s_i) = (s_i - \bar s)/\hat\sigma_s$ with the sample mean and the sample
standard deviation over all $N$ agents. Status is standardized on its own
lognormal scale, so $z(s)$ is right-skewed. Agent $i$ also carries three taste
multipliers $m^{\alpha}_i, m^{\beta}_i, m^{\gamma}_i$, independent lognormal with
log-mean $-\sigma_\tau^2/2$ and log-sd $\sigma_\tau$, so each has mean one.

The utility of $i$ for a candidate $j$ on the other side is
\begin{equation}
U_{ij} = m^{\alpha}_i \alpha\, z(s_j) + m^{\beta}_i \beta\, z(a_j)
       - m^{\gamma}_i \gamma\, d_{ij}
       + \lambda\, \mathbf{1}[i \in W]\, \mathbf{1}\!\left[z(s_j) \ge z(s_i) + \delta\right]
       + \sigma_L L_{ij},
\label{eq:utility}
\end{equation}
where $d_{ij} = \lVert z(x_i) - z(x_j) \rVert_2$ is the distance between the
standardized trait vectors. The fourth term is the status-surplus bonus: an
acceptor gains $\lambda$ from a candidate whose standardized status exceeds her
own by at least $\delta$. It is the only asymmetry between the sides, and it is
one variant of the utility rather than the subject of the paper
(\S\ref{sec:where}). The last term is persistent pairwise chemistry: $L_{ij}$ is a
standard logistic draw computed from a hash of the run seed and two per-agent
keys, fixed for the life of the pair and asymmetric ($L_{ij} \ne L_{ji}$). Its
standard deviation is $\sigma_L \pi/\sqrt3$, which at the shipped $\sigma_L = 1.5$
is $2.72$ utility units.

Agents also carry an intent $\tau_i \sim \mathrm{Beta}(2,2)$, a values score $g_i
\sim \mathrm{Beta}(2,2)$ and an eligibility flag for the institutional comparison.
They enter only that comparison (\S\ref{sec:institutional},
Appendix~\ref{app:hypotheses}) and change no exposure quantity.

\subsection{Acceptance thresholds and the gate}
\label{sec:gate}

Each agent has an acceptance threshold $\theta_i$ in utility units, and its
\emph{admissible set} is $A_i = \{ j : U_{ij} \ge \theta_i \}$. At round 0,
$\theta_i$ is the $n_A$-th largest of $i$'s utilities over the other side, with
$n_A = \mathrm{round}(\phi N/2)$, so every agent starts with exactly $n_A$
admissible candidates. After each round,
\begin{equation}
\theta_i \leftarrow \theta_i + \eta \left( \frac{m_i}{p_i} - \theta_i \right)
\quad \text{if } p_i > 0,
\label{eq:theta}
\end{equation}
where $p_i$ is the number of proposals $i$ made (on $M$) or received (on $W$) in
that round's matching, $m_i \in \{0, 1\}$ is whether $i$ was matched, and $\eta =
0.1$. An agent with no proposal keeps its threshold. The yield $m_i/p_i$ lies in
$[0, 1]$ while the initial threshold is the tenth-highest utility an agent has,
far above one, so the update lowers $\theta_i$ geometrically toward the unit
interval for every agent who takes part in the matching.

The threshold can change what an agent sees only if the admissible set is
shorter than the shortlist. The candidate set is the other side, of size $N/2$,
so $|A_i| = \phi N/2$ at round 0 and the gate binds then exactly when
\begin{equation}
\phi < \frac{2k}{N}.
\label{eq:gate}
\end{equation}
At $N = 1000$ and $k = 20$ this is $\phi < 0.04$; the shipped $\phi = 0.02$ gives
ten admissible candidates against a shortlist of twenty. As thresholds fall the
admissible sets grow past $k$ and the gate stops binding, within a few rounds
for agents who propose or are proposed to.

\subsection{The training log}
\label{sec:log}

Each round the ranker trains on a log built afresh from the current population
and thresholds. For every agent $i$ on either side:
\begin{itemize}
\item the \emph{positives} are the $\min(k, |A_i|)$ members of $A_i$ with the
  highest $U_{ij}$, each labeled 1;
\item the \emph{negatives} are $k n_{-}$ candidates drawn uniformly without
  replacement from the rest of the other side, each labeled 0, with $n_{-} = 4$.
\end{itemize}
A positive is determined by the utility and the threshold; there is no
randomness in it. The negative draw for agent $i$ in round $r$ comes from a
generator seeded with (run seed, $i$, $k$, $r$), so every round has fresh
negatives and the log does not depend on the thread count. Logs are not
accumulated: the ranker sees one round's log at a time, and what it retains is in
its weights. At the shipped values an agent contributes at most $k(1 + n_{-}) =
100$ rows to a log, and \rowsPerRound{} on average over a run.

\textbf{The log does not depend on the ranker's shortlists.} Its positives are
exactly the shortlists an oracle with the same thresholds would hand out, and its
negatives are uniform. The ranker's own choices reach its training data through
one channel only: they decide who proposes to whom, which moves thresholds by
\eqref{eq:theta}, which changes admissible sets. This is the fact that answers
the convergence objection of \S\ref{sec:intro}, and \S\ref{sec:what-ends} tests
it. It also means that, once the gate stops binding, the positives are the
oracle's top-$k$ lists by true utility: the ranker imitates the benchmark that the
gap is measured against, and the gap is the rate at which that imitation
converges. It is not the exposure of a ranker trained on engagement that its own
lists produce.

\subsection{The ranker}
\label{sec:ranker}

The ranker is neural collaborative filtering in the fused form of
\citet{he2017}: a generalized matrix factorization branch and a multilayer
perceptron branch over one id space, in which every agent is both a user and a
candidate. Each branch has its own user and item embeddings of dimension 64. The
perceptron maps the concatenated pair embedding through layers of width 128, 64
and 32 with ReLU activations, and a linear output layer reads the element-wise
product of the factorization embeddings beside the perceptron's output to give a
logit $\hat y_{uv}$. Embeddings are initialized $N(0, 0.01^2)$ and dense layers
with the library default. Training minimizes binary cross-entropy on the logits
with positive weight 4, using Adam \citep{kingma2015} at learning rate $10^{-3}$
and weight decay $10^{-5}$, in batches of 4096 rows shuffled by a seeded
generator. The weights and the optimizer state persist across rounds. One
training pass over a round's log is about \stepsPerPass{} optimizer steps at the
shipped values; the dose experiment of \S\ref{sec:what-ends} changes the number
of passes.

Agent $u$'s shortlist $S_u$ is the $k$ candidates on the other side with the
highest $\hat y_{uv}$. It is \emph{not} restricted to $A_u$: the ranker ranks the
whole other side. An agent's \emph{exposure} $e_j$ in a round is the number of
shortlists that contain it, $0 \le e_j \le N/2$.

The architecture is deliberately standard. \citet{rendle2020} show that a tuned
dot product beats this learned similarity at recommendation accuracy; no result
here depends on the ranker being a good one, only on its being learned from the
log above, and \S\ref{sec:limits} returns to what a different ranker could
change.

\subsection{Clearing}
\label{sec:clearing}

The matching is deferred acceptance \citep{gale1962} with $M$ proposing, run on
the shortlists: a proposer works down his own shortlist, and an acceptor holds
the best proposal from among the proposers \emph{on her shortlist} and rejects
everyone else, including proposers she has never been shown. The result is
stable with respect to the shortlists and in general not with respect to true
preferences (Table~\ref{tab:validation}). With $M$ proposing it is the
proposer-optimal stable matching of the truncated lists and the acceptor-pessimal
one \citep{roth1982,knuth1996}. Every arm uses the same proposing side, so the
comparisons below are between rankers, not between directions.

\subsection{Comparison arms}
\label{sec:arms}

Four rules produce shortlists on the same population, utilities, seed and
matching code; only the rule differs, and each arm runs its own threshold
dynamics.
\begin{itemize}
\item \textbf{NCF}: the learned ranker above.
\item \textbf{Oracle}: the top $k$ of $A_i$ by true utility, shorter when
  $|A_i| < k$. It is the benchmark a learned ranker falls short of, not an
  achievable institution: it knows the chemistry term, which nothing observable
  reveals.
\item \textbf{Random}: $\min(k, |A_i|)$ members of $A_i$ drawn uniformly without
  replacement, from a generator seeded with (run seed, round).
\item \textbf{Popularity}: one list per side, the $k$ agents of the other side
  with the highest exposure in the previous round, ties to the lower index. In
  round 0, which has no history, it ranks by true utility.
\end{itemize}

\subsection{Metrics}
\label{sec:metrics}

All metrics are per round. With exposures $e_1, \ldots, e_N$ over all agents of
both sides, sorted ascending as $e_{(1)} \le \cdots \le e_{(N)}$:
\begin{align}
G &= \frac{2 \sum_{i=1}^{N} i\, e_{(i)}}{N \sum_{i=1}^{N} e_{(i)}} - \frac{N+1}{N}
  && \text{exposure Gini,} \label{eq:gini}\\
R &= \frac{1}{N} \left| \{ j : e_j > 0 \} \right|
  && \text{reach,} \notag\\
\Delta &= G_{\mathrm{NCF}} - G_{\mathrm{oracle}}
  && \text{the gap, same seed and round.} \notag
\end{align}
The gap is the paper's quantity. It measures concentration in excess of the
oracle's, and the oracle's own Gini is near $\giniOracleTail$
(\S\ref{sec:transient}): that much concentration comes from preferences, because
agents with high status and attractiveness are on many agents' true top-$k$. A
positive gap is concentration the ranker adds to that, not concentration as
such.

Two more metrics are reported. The \emph{clearing rate} is $2 \mu / N$ for $\mu$
matched pairs. \emph{Algorithmic loss} is $(W^{*} - W)/W^{*}$, where $W = \sum
(U_{ij} + U_{ji})$ over matched pairs and $W^{*}$ is the same sum for deferred
acceptance on complete true-preference lists. In the fully concentrated market,
where every list on a side is the same $k$ agents, exactly $k$ pairs can form and
the clearing rate is $2k/N$; when lists spread it is larger (Supplementary note
S1).

\subsection{An open market}
\label{sec:open}

Every run except those of \S\ref{sec:arrivals} has a closed population: no agent
arrives or leaves. In the open-market runs a share $\nu$ of the $N$ slots changes
hands at the start of every round after the first. The departing agents are drawn
uniformly at random, so tenure is geometric with mean $1/\nu$ rounds, and each is
replaced in the same slot by a fresh agent: new traits, tastes, chemistry key,
intent and values from the distributions of \S\ref{sec:utility}; a threshold
calibrated against the market it enters; the ranker's four embedding rows for the
slot re-drawn from the initial distribution, with the shared layers kept. This is
replacement, not net arrival: the market's size is fixed. The departures and the
newcomers depend only on the run seed and the round, so every arm on a seed sees
the same population at every round.

The code also supports generational turnover, in which the least-matched fifth
of the population is replaced at fixed intervals. No experiment reported here
uses it.

\subsection{The second environment}
\label{sec:env-b}

The institutional comparison of \S\ref{sec:institutional} sets the market above,
Environment A, against Environment B: a subgraph of the same population in which
every member sees every other member. An agent is eligible for B if its
eligibility flag is set and $g_i \ge g_{\min} = 0.3$; a share $\phi_B$ of each
side is drawn into B, and the sides are balanced by dropping the excess. Inside B
the shortlist is the oracle rule restricted to members, and proposing costs
$c_B$. An agent $i$ proposes to $j$ only if
\begin{equation}
\tau_i\, U_{ij}\, p > c_B,
\label{eq:signal}
\end{equation}
where $p$ is the acceptance rate measured on the previous round's log for $i$'s
side (0.5 before any measurement). A matched pair survives the round with
probability
\begin{equation}
\sigma_{ij} = \sigma_0 (\tau_i \tau_j)^{\sigma_\alpha} g_{ij}^{\sigma_\beta},
\qquad g_{ij} = 1 - |g_i - g_j|,
\label{eq:survival}
\end{equation}
with $\sigma_0 = 1$ and $\sigma_\alpha = \sigma_\beta = 0.5$. When a pair dissolves
the member with lower intent defects; the partner loses $L(1 - \tau_{\mathrm{def}})$
with $L = 1$, and in B the defector is charged $d = 0.5$. In A the same survival
draw applies with $d = 0$ and no proposal cost. Welfare in a round is $\sum (\tau_i
U_{ij} + \tau_j U_{ji})$ over matched pairs, less proposal costs, less dissolution
losses. Agents B excludes are scored at the payoff they would get in A with B's
matched agents removed from A's market, so that excluding an agent cannot raise
B's welfare by itself.

\begin{algorithm}[t]
\caption{One round $r$ of the market. $\mathcal{L}$ is the round's training log.}
\label{alg:round}
\begin{algorithmic}[1]
\If{$r > 0$ and $\nu > 0$} replace a share $\nu$ of the slots (\S\ref{sec:open})
\EndIf
\State $\mathcal{L} \gets$ positives and negatives for every agent (\S\ref{sec:log})
\State take $\rho$ training passes of the ranker over $\mathcal{L}$ ($\rho = 1$ shipped)
\For{each agent $u$}
  \State $S_u \gets$ the top $k$ of the other side by score
\EndFor
\State match by deferred acceptance, $M$ proposing, on $\{S_u\}$ (\S\ref{sec:clearing})
\State record exposure $e$, $G$, $R$, the clearing rate and welfare
\For{each agent $i$ with $p_i > 0$}
  \State $\theta_i \gets \theta_i + \eta\,(m_i/p_i - \theta_i)$
\EndFor
\end{algorithmic}
\end{algorithm}

\begin{table}[t]
\centering
\small
\begin{tabular}{llll}
\hline
Symbol & Meaning & Shipped & Swept \\
\hline
$N$ & agents & 1000 & 500, 2000 \\
$k$ & shortlist length & 20 & not swept \\
$\phi$ & admissible share at round 0 & 0.02 & $20/N$ with $N$; 0.2 (inert gate) \\
$\eta$ & threshold step & 0.1 & not swept \\
$\nu$ & replacement share per round & 0 & 0.005, 0.01, 0.02 \\
\hline
$\alpha$ & status weight & 1 & not swept \\
$\beta$ & attractiveness weight & 2 & 1, 3 \\
$\gamma$ & distance penalty & 0.3 & not swept \\
$\lambda$ & status-surplus bonus & 1.5 & not swept \\
$\delta$ & status margin (sd) & 0.5 & 0, 1 \\
$\sigma_L$ & chemistry scale & 1.5 & 0.5, 3.0 \\
$\sigma_\tau$ & taste spread (log-sd) & 0.5 & 0.2, 1.0 \\
\hline
$s$ & status, lognormal (log-mean, log-sd) & 0, 0.8 & not swept \\
$a$ & attractiveness, normal, clipped & 0.5, 0.15 & not swept \\
$c$ & sociability, normal, clipped & 0.5, 0.15 & not swept \\
$t$ & age, normal & 30, 6 & not swept \\
$\tau, g$ & intent, values: Beta & (2, 2) & not swept \\
\hline
 & embedding dimension & 64 & not swept \\
 & perceptron widths & 128, 64, 32 & not swept \\
 & learning rate, weight decay & $10^{-3}$, $10^{-5}$ & not swept \\
 & positive weight, negatives per positive & 4, 4 & not swept \\
$\rho$ & training passes per round & 1 & 0.5, 2, 4 \\
 & exploration share & 0 & 0.05, 0.2 \\
 & warm-start passes & 0 & 50 \\
\hline
$\phi_B$ & subgraph share & 0.5 & 0.02 to 1.00 (ten levels) \\
$c_B$ & proposal cost & 0.25 & 0.01, 0.05, 0.10, 0.25 \\
$g_{\min}$ & values filter & 0.3 & not swept \\
$d$ & defection penalty & 0.5 & not swept \\
$\sigma_0, \sigma_\alpha, \sigma_\beta$ & survival & 1, 0.5, 0.5 & not swept \\
$L$ & partner loss & 1 & not swept \\
\hline
\end{tabular}

\caption{Every parameter of the model, with its value in the shipped
configuration and the values swept. All are modeling choices: none is estimated
from data. The swept values are those of the sensitivity sweep
(\S\ref{sec:where}), the dose, exploration and population-size experiments
(\S\ref{sec:what-ends}) and the open market (\S\ref{sec:arrivals}).}
\label{tab:params}
\end{table}
